\documentclass[a4paper,landscape,8pt]{extarticle}
\usepackage[landscape,margin=0.7cm,top=0.5cm,bottom=0.7cm]{geometry}
\usepackage[utf8]{inputenc}
\usepackage[T1]{fontenc}
\usepackage[ngerman]{babel}
\usepackage{helvet}
\renewcommand{\familydefault}{\sfdefault}
\usepackage{xcolor}
\usepackage{tikz}
\usetikzlibrary{positioning, shapes.geometric, calc}
\usepackage{enumitem}
\usepackage{multicol}
\usepackage{array}
\usepackage{fancyhdr}
\usepackage{amssymb}

% Farben
\definecolor{ipadsilver}{HTML}{8E8E93}
\definecolor{splitblue}{HTML}{007AFF}
\definecolor{pencilgold}{HTML}{F5A623}
\definecolor{appgreen}{HTML}{34C759}
\definecolor{lightgray}{HTML}{F5F5F5}
\definecolor{bordergray}{HTML}{BBBBBB}
\definecolor{warnred}{HTML}{C62828}

\pagestyle{fancy}
\fancyhf{}
\renewcommand{\headrulewidth}{0pt}
\fancyfoot[C]{\footnotesize Seite \thepage{} von 3}

\setlength{\parindent}{0pt}
\setlength{\parskip}{2pt}

% Kompakte Listen
\setlist{nosep, leftmargin=1em, itemsep=0pt, topsep=1pt}

\begin{document}

% ============================================================================
% SEITE 1: INFORMATIONSBLATT
% ============================================================================

\noindent{\Large\bfseries Das iPad kennenlernen: Mehr Platz, mehr Möglichkeiten} \hfill
{\footnotesize\textbf{Lernziele:} Unterschiede zum iPhone | Dock nutzen | Multitasking verstehen}

\vspace{1mm}
\hrule
\vspace{1mm}

% Drei-Spalten-Layout
\begin{minipage}[t]{0.32\linewidth}
\begin{center}
\colorbox{ipadsilver!15}{%
\begin{minipage}{0.95\linewidth}
\vspace{2mm}
\begin{center}
{\Large\bfseries\textcolor{ipadsilver}{iPad vs. iPhone}}\\[1pt]
{\footnotesize\itshape Was ist anders?}
\end{center}
\vspace{2mm}
\end{minipage}%
}
\end{center}

\vspace{1mm}

\textbf{Was ist ein iPad?} Ein \textbf{großes iPhone} -- gleiche Bedienung, mehr Platz.

\textbf{Analogie:} iPhone = \textbf{Notizblock}, iPad = \textbf{Schreibtisch} -- gleiche Arbeit, aber bequemer.

\textbf{Was ist gleich?}
\begin{itemize}
\item Gleiche Apps (aus dem App Store)
\item Gleiche Apple-ID und iCloud
\item Gleiche Gesten (Tippen, Wischen...)
\item Gleiche Einstellungen-App
\end{itemize}

\textbf{Was ist anders?}
\begin{itemize}
\item \textbf{Größerer Bildschirm} -- mehr sehen
\item \textbf{Dock} zeigt mehr Apps (bis zu 15)
\item \textbf{Zwei Apps gleichzeitig} möglich
\item \textbf{Apple Pencil} zum Schreiben
\item Kein Telefonieren (meist)
\end{itemize}

\textbf{Brauche ich beides?}\\
Nicht unbedingt. Viele nutzen das iPad zu Hause und das iPhone unterwegs.

\end{minipage}%
\hfill%
\begin{minipage}[t]{0.32\linewidth}
\begin{center}
\colorbox{splitblue!10}{%
\begin{minipage}{0.95\linewidth}
\vspace{2mm}
\begin{center}
{\Large\bfseries\textcolor{splitblue}{Das Dock}}\\[1pt]
{\footnotesize\itshape Deine Werkzeugleiste}
\end{center}
\vspace{2mm}
\end{minipage}%
}
\end{center}

\vspace{1mm}

\textbf{Was ist das Dock?} Die \textbf{Leiste unten} mit deinen wichtigsten Apps.

\textbf{Analogie:} Wie die \textbf{Werkzeugleiste} in der Werkstatt -- die meistgenutzten Werkzeuge griffbereit.

\textbf{Besonderheit beim iPad:}
\begin{itemize}
\item Mehr Platz als beim iPhone
\item \textbf{Immer erreichbar} -- auch in Apps!
\item Von unten nach oben wischen = Dock erscheint
\end{itemize}

\textbf{Dock anpassen:}
\begin{enumerate}
\item App auf Home-Bildschirm \textbf{lange drücken}
\item Wenn sie wackelt: ins Dock \textbf{ziehen}
\item Außerhalb tippen = fertig
\end{enumerate}

\textbf{Zwei Bereiche im Dock:}
\begin{itemize}
\item \textbf{Links:} Deine festen Apps
\item \textbf{Rechts} (nach Strich): Zuletzt benutzte Apps
\end{itemize}

\end{minipage}%
\hfill%
\begin{minipage}[t]{0.32\linewidth}
\begin{center}
\colorbox{appgreen!10}{%
\begin{minipage}{0.95\linewidth}
\vspace{2mm}
\begin{center}
{\Large\bfseries\textcolor{appgreen}{Home-Bildschirm}}\\[1pt]
{\footnotesize\itshape Mehr Platz, mehr Apps}
\end{center}
\vspace{2mm}
\end{minipage}%
}
\end{center}

\vspace{1mm}

\textbf{Mehr Apps pro Seite:}\\
Durch den größeren Bildschirm passen \textbf{mehr App-Symbole} auf eine Seite.

\textbf{Widgets:} Kleine Info-Kästchen auf dem Home-Bildschirm (Wetter, Kalender, Uhr...).

\textbf{Widget hinzufügen:}
\begin{enumerate}
\item Lange auf freie Stelle drücken
\item Oben links auf \textbf{+} tippen
\item Widget auswählen und platzieren
\end{enumerate}

\textbf{App-Mediathek:}\\
Ganz nach links wischen -- alle Apps alphabetisch sortiert.

\vspace{2mm}

\textbf{Gesten wie beim iPhone:}

\begin{tabular}{@{}p{3cm}p{4.5cm}@{}}
Nach oben wischen & Home-Bildschirm \\
Oben rechts wischen & Kontrollzentrum \\
Mitte runter wischen & Suche \\
\end{tabular}

\end{minipage}

\vspace{1mm}

% Zusammenfassung
\begin{center}
\fcolorbox{bordergray}{lightgray}{%
\begin{minipage}{0.98\linewidth}
\centering
\vspace{1.5mm}
\textbf{Merke:} Das iPad funktioniert wie ein \textbf{großes iPhone}. Die Bedienung ist gleich -- aber du hast \textbf{mehr Platz} und kannst \textbf{zwei Apps gleichzeitig} nutzen.
\vspace{1.5mm}
\end{minipage}%
}
\end{center}

\newpage

% ============================================================================
% SEITE 2: MULTITASKING & APPLE PENCIL
% ============================================================================

\begin{center}
{\LARGE\bfseries Multitasking \& Apple Pencil}\\[3pt]
{\normalsize Zwei Apps gleichzeitig nutzen und mit dem Stift schreiben}
\end{center}

\vspace{1mm}
\hrule
\vspace{2mm}

\begin{multicols}{2}

\section*{\textcolor{splitblue}{Split View -- Zwei Apps nebeneinander}}

\textbf{Was ist das?} Du kannst \textbf{zwei Apps gleichzeitig} auf dem Bildschirm sehen und nutzen.

\textbf{Analogie:} Wie zwei Bücher \textbf{nebeneinander} auf dem Tisch -- du kannst in beiden lesen.

\textbf{Wofür ist das gut?}
\begin{itemize}
\item E-Mail lesen und gleichzeitig Notizen machen
\item Rezept anschauen und Einkaufsliste schreiben
\item Foto ansehen und in Nachricht einfügen
\end{itemize}

\textbf{So aktivierst du Split View:}
\begin{enumerate}
\item Öffne die \textbf{erste App}
\item Wische von unten, um das \textbf{Dock} zu zeigen
\item Ziehe eine \textbf{zweite App} aus dem Dock an den Rand
\item Loslassen -- fertig!
\end{enumerate}

\textbf{Größe anpassen:}\\
Den \textbf{grauen Balken} in der Mitte nach links oder rechts ziehen.

\textbf{Split View beenden:}\\
Grauen Balken ganz nach links oder rechts ziehen.

\vspace{3mm}

\section*{\textcolor{ipadsilver}{Slide Over -- App im Fenster}}

\textbf{Was ist das?} Eine kleine App \textbf{schwebt über} der großen App.

\textbf{So geht's:}
\begin{enumerate}
\item App aus dem Dock in die \textbf{Mitte} ziehen (nicht an den Rand)
\item Die App erscheint als schwebendes Fenster
\end{enumerate}

\textbf{Slide Over verschieben:}\\
Am grauen Strich oben anfassen und ziehen.

\textbf{Slide Over ausblenden:}\\
Nach rechts vom Bildschirm wischen -- kommt zurück, wenn du vom rechten Rand wischt.

\columnbreak

\section*{\textcolor{pencilgold}{Apple Pencil -- Der Stift fürs iPad}}

\textbf{Was ist das?} Ein spezieller \textbf{Stift von Apple} zum Schreiben und Zeichnen auf dem iPad.

\textbf{Analogie:} Wie ein \textbf{Kugelschreiber} -- aber für den Bildschirm.

\textbf{Wichtig:} Nicht jeder Stift funktioniert! Es muss ein \textbf{Apple Pencil} sein.

\vspace{2mm}

\textbf{Was kann man damit machen?}
\begin{itemize}
\item \textbf{Handschriftlich schreiben} (wird in Text umgewandelt!)
\item \textbf{Unterschreiben} in Dokumenten
\item \textbf{Zeichnen} und Malen
\item \textbf{Markieren} in PDFs und Fotos
\item \textbf{Notizen} wie auf Papier
\end{itemize}

\vspace{2mm}

\textbf{Apple Pencil laden:}
\begin{itemize}
\item \textbf{1. Generation:} In Lightning-Anschluss stecken
\item \textbf{2. Generation:} An die Seite des iPads legen (magnetisch)
\end{itemize}

\vspace{2mm}

\textbf{Kritzeln (Scribble):}\\
Mit dem Pencil in \textbf{jedes Textfeld} schreiben -- das iPad wandelt deine Handschrift in getippten Text um!

\vspace{3mm}

\begin{center}
\fcolorbox{warnred}{warnred!8}{%
\begin{minipage}{0.9\linewidth}
\vspace{2mm}
\textbf{Kein Apple Pencil?}\\[2pt]
Kein Problem! Das iPad funktioniert auch ohne Stift -- alles geht auch mit dem Finger. Der Pencil ist \textbf{optional}.
\vspace{2mm}
\end{minipage}}
\end{center}

\vspace{3mm}

\begin{center}
\fcolorbox{splitblue}{splitblue!8}{%
\begin{minipage}{0.9\linewidth}
\vspace{2mm}
\textbf{Tipp: Schnelle Notiz}\\[2pt]
Mit dem Apple Pencil in die \textbf{rechte untere Ecke} tippen -- sofort öffnet sich eine neue Notiz!
\vspace{2mm}
\end{minipage}}
\end{center}

\end{multicols}

\newpage

% ============================================================================
% SEITE 3: ÜBUNGEN & CHECKLISTE
% ============================================================================

\begin{center}
{\LARGE\bfseries Übungen \& Checkliste}\\[3pt]
{\normalsize Teste dein Wissen und übe die iPad-Funktionen}
\end{center}

\vspace{1mm}
\hrule
\vspace{2mm}

\begin{multicols}{2}

\textbf{\large Teil A: iPad vs. iPhone}

Was stimmt? Kreuze an: R = richtig, F = falsch

\vspace{2mm}

\renewcommand{\arraystretch}{1.4}
\begin{tabular}{@{}p{8cm}cc@{}}
& R & F \\
1. Das iPad hat die gleichen Apps wie das iPhone. & $\square$ & $\square$ \\
2. Auf dem iPad kann ich zwei Apps gleichzeitig nutzen. & $\square$ & $\square$ \\
3. Das iPad braucht eine andere Apple-ID als das iPhone. & $\square$ & $\square$ \\
4. Jeder Stift funktioniert auf dem iPad. & $\square$ & $\square$ \\
5. Das Dock beim iPad zeigt mehr Apps als beim iPhone. & $\square$ & $\square$ \\
\end{tabular}
\renewcommand{\arraystretch}{1}

\textit{\footnotesize Lösungen: R, R, F, F, R}

\vspace{4mm}

\textbf{\large Teil B: Praktische Übungen}

Führe diese Aktionen auf deinem iPad aus und hake ab:

\vspace{2mm}

\begin{itemize}[label=$\square$, itemsep=3pt]
\item Öffne das \textbf{Dock} (von unten nach oben wischen)
\item Öffne das \textbf{Kontrollzentrum} (von oben rechts wischen)
\item Öffne die \textbf{Suche} (von der Mitte nach unten wischen)
\item Füge eine App zum \textbf{Dock} hinzu
\item Öffne \textbf{zwei Apps nebeneinander} (Split View)
\item Ändere die \textbf{Größe} der beiden Apps
\item \textbf{Beende} Split View
\end{itemize}

\vspace{4mm}

\textbf{\large Teil C: Was machst du?}

\vspace{2mm}

\textbf{Situation 1:} Du willst eine E-Mail lesen und gleichzeitig etwas nachschlagen.

\rule{7cm}{0.4pt}

\vspace{3mm}

\textbf{Situation 2:} Du suchst eine App, findest sie aber nicht auf dem Home-Bildschirm.

\rule{7cm}{0.4pt}

\columnbreak

\textbf{\large Checkliste: Das kann ich jetzt}

\vspace{2mm}

\begin{center}
\fcolorbox{ipadsilver}{ipadsilver!10}{%
\begin{minipage}{0.95\linewidth}
\vspace{2mm}
\textbf{\textcolor{ipadsilver}{iPad-Grundlagen}}
\begin{itemize}[label=$\square$, itemsep=2pt]
\item Ich verstehe, dass das iPad wie ein \textbf{großes iPhone} funktioniert
\item Ich weiß, was das \textbf{Dock} ist und wie ich es öffne
\item Ich kann Apps zum Dock \textbf{hinzufügen}
\item Ich kann die \textbf{App-Mediathek} öffnen
\end{itemize}
\vspace{2mm}
\end{minipage}}
\end{center}

\vspace{2mm}

\begin{center}
\fcolorbox{splitblue}{splitblue!10}{%
\begin{minipage}{0.95\linewidth}
\vspace{2mm}
\textbf{\textcolor{splitblue}{Multitasking}}
\begin{itemize}[label=$\square$, itemsep=2pt]
\item Ich kann \textbf{zwei Apps nebeneinander} öffnen (Split View)
\item Ich kann die \textbf{Größe} der Apps anpassen
\item Ich kann Split View \textbf{beenden}
\item Ich weiß, was \textbf{Slide Over} ist
\end{itemize}
\vspace{2mm}
\end{minipage}}
\end{center}

\vspace{2mm}

\begin{center}
\fcolorbox{pencilgold}{pencilgold!10}{%
\begin{minipage}{0.95\linewidth}
\vspace{2mm}
\textbf{\textcolor{pencilgold}{Apple Pencil (falls vorhanden)}}
\begin{itemize}[label=$\square$, itemsep=2pt]
\item Ich weiß, wie ich den Pencil \textbf{lade}
\item Ich weiß, dass ich mit \textbf{Kritzeln} handschriftlich in Textfelder schreiben kann
\item Ich kenne den Trick mit der \textbf{schnellen Notiz}
\end{itemize}
\vspace{2mm}
\end{minipage}}
\end{center}

\vspace{3mm}

\begin{center}
\fcolorbox{bordergray}{white}{%
\begin{minipage}{0.95\linewidth}
\vspace{2mm}
\textbf{Meine Notizen}\\[2pt]
Mein iPad-Modell: \dotfill\\[2mm]
Habe ich einen Apple Pencil? $\square$ Ja \quad $\square$ Nein\\[2mm]
Was ich noch üben möchte: \dotfill
\vspace{2mm}
\end{minipage}}
\end{center}

\end{multicols}

\vfill

\begin{center}
\footnotesize\textit{Stand: \today{} -- Das iPad ist wie ein großes iPhone. Probiere Split View aus -- es ist sehr praktisch!}
\end{center}

\end{document}
